\chapter{Project Requirements}
\label{chap:requirements}

This chapter outlines the requirements for the successful completion of Vital
Guardian, organised into three categories: the Event Response Table (Use Cases),
Functional Requirements per module, and Non-Functional Requirements.

% =============================================================
\section{Use-Case and Event Response Table}
\label{sec:event_response}

The system is designed to detect specific environmental stimuli and trigger
calculated responses. Table \ref{tab:event_response} details the full range of
critical and non-critical events the system monitors, and the corresponding
responses determined by the Cognitive Core.

\begin{table}[H]
    \centering
    \caption{Event Response Table for Vital Guardian (Part 1 of 2)}
    \label{tab:event_response}
    \renewcommand{\arraystretch}{1.4}
    \begin{tabular}{|p{3.8cm}|p{2cm}|p{5.2cm}|p{1.8cm}|}
    \hline
    \textbf{Event (Stimulus)} & \textbf{Source} & \textbf{System Response} & \textbf{Priority} \\ \hline

    Patient experiences a sudden rapid descent from the bed, followed by a gasp. &
    Vision and Audio &
    Cognitive Core fuses the fall event and the Tier-3 distress audio, confirms
    the emergency, and generates a natural language alert: ``Patient has fallen
    and is in distress.'' &
    CRITICAL \\ \hline

    Erratic, seizure-like full-body movements detected, with wheezing or
    laboured breathing audible. &
    Vision and Audio &
    MoViNet-A2 and YAMNet both fire simultaneously. Cognitive Core correlates
    dangerous movement with respiratory distress and flags a medical emergency. &
    CRITICAL \\ \hline

    Patient screams or calls for help. Keyword ``Help'' or ``Dard'' detected in
    the audio stream. &
    Audio &
    Faster-Whisper transcribes the keyword. Cognitive Core generates an
    immediate alert containing the detected phrase and timestamp. &
    CRITICAL \\ \hline

    Patient rolls to the edge of the bed and moaning is detected. &
    Vision and Audio &
    MediaPipe detects the hip landmark approaching the bed boundary.
    YAMNet detects moaning. Cognitive Core generates alert:
    ``Patient showing fall risk near bed edge.'' &
    HIGH \\ \hline

    Patient fall detected with no accompanying audio distress signal. &
    Vision &
    MoViNet-A2 fires on the fall clip. Cognitive Core performs Tier-2
    verification on the vision event alone and registers a high-priority alert
    to ensure patient safety. &
    HIGH \\ \hline

    Patient detected attempting to remove medical equipment, with rapid
    irregular arm movements near the torso. &
    Vision &
    MediaPipe detects sustained wrist proximity to chest region. Cognitive Core
    classifies the event as a high-risk interference attempt and alerts staff. &
    HIGH \\ \hline

    \multicolumn{4}{r}{\small\textit{Continued on next page\ldots}} \\
    \end{tabular}
\end{table}

\begin{table}[H]
    \centering
    \caption{Event Response Table for Vital Guardian (Part 2 of 2)}
    \label{tab:event_response_2}
    \renewcommand{\arraystretch}{1.4}
    \begin{tabular}{|p{3.8cm}|p{2cm}|p{5.2cm}|p{1.8cm}|}
    \hline
    \textbf{Event (Stimulus)} & \textbf{Source} & \textbf{System Response} & \textbf{Priority} \\ \hline

    Prolonged period of complete patient stillness detected. No movement for
    an extended duration beyond the configured threshold. &
    Vision &
    Visual Guardian detects the absence of expected baseline micro-movement.
    Cognitive Core generates a medium-priority welfare check alert. &
    MEDIUM \\ \hline

    Coughing or choking sounds detected persistently across multiple consecutive
    audio chunks. &
    Audio &
    YAMNet classifies repeated Tier-1 audio events. Cognitive Core escalates
    to MEDIUM priority after sustained duration and alerts nursing staff. &
    MEDIUM \\ \hline

    Visitor conversation detected. Speech activity continues beyond the
    configured visitor duration threshold. &
    Audio &
    Silero VAD activates Visitor Mode. Speech transcription and keyword
    spotting are automatically paused. No alert is generated. Dashboard
    displays ``Visitor Mode Active.'' &
    LOW \\ \hline

    Patient adjusts sleeping position calmly. No audio distress signal
    is detected. &
    Vision &
    MediaPipe detects movement but torso angle and kinematic variance remain
    within safe bounds. Cognitive Core logs the event as routine. No alarm
    generated to reduce alarm fatigue. &
    LOW \\ \hline

    Loud environmental noise detected (equipment alarm, door slam) with no
    corresponding visual patient event. &
    Audio &
    YAMNet classifies the sound. Cognitive Core finds no corroborating vision
    event and suppresses the alert, logging it as a false positive candidate. &
    LOW \\ \hline

    Patient is outside the camera frame. No bounding box returned by the
    person detector for a sustained period. &
    Vision &
    YOLOv11n fails to return a detection for the configured number of
    consecutive frames. System generates an alert: ``Patient not visible.
    Manual check required.'' &
    HIGH \\ \hline

    \end{tabular}
\end{table}


% =============================================================
\section{Functional Requirements}
\label{sec:functional_reqs}

This section describes the functional requirements organised by module.

\subsection{Module 1: The Visual Guardian}
\begin{enumerate}
    \item \textbf{FR-V-1:} The system shall detect and track the patient within
    the camera's field of view in real-time using a YOLOv11n object detector.
    \item \textbf{FR-V-2:} The system shall detect when a patient has fallen from
    the bed using a MoViNet-A2 spatiotemporal classifier.
    \item \textbf{FR-V-3:} The system shall detect dangerous or unusual movements
    (e.g., seizures) using a dedicated temporal action classifier and kinematic
    landmark analysis.
    \item \textbf{FR-V-4:} The system shall provide a configurable option to
    pixelate the patient's face in the live video feed to protect privacy.
\end{enumerate}

\subsection{Module 2: The Auditory Watchdog}
\begin{enumerate}
    \item \textbf{FR-A-1:} The system shall classify non-verbal distress sounds,
    including gasps, moans, and severe coughing, using the YAMNet audio classifier.
    \item \textbf{FR-A-2:} The system shall detect predefined keywords in English
    (``Help'', ``Pain'') and Urdu (``Madad'', ``Dard'') using an offline
    speech-to-text transcription model.
    \item \textbf{FR-A-3:} The system shall implement a Privacy Shield by using
    Voice Activity Detection to identify continuous speech (indicating a visitor
    conversation) and automatically pause speech transcription.
\end{enumerate}

\subsection{Module 3: The Cognitive Core}
\begin{enumerate}
    \item \textbf{FR-C-1:} The system shall fuse structured event data from the
    Visual and Auditory modules to perform multi-modal contextual reasoning.
    \item \textbf{FR-C-2:} The system shall generate clear, human-readable alert
    messages in natural language suitable for immediate comprehension by nursing
    staff, without requiring any specialised training.
\end{enumerate}

% =============================================================
\section{Non-Functional Requirements}
\label{sec:nfr}

\subsection{Reliability}
\begin{enumerate}
    \item \textbf{REL-1:} The system shall provide robust runtime stability through
    Docker containerisation, allowing the FastAPI application and PostgreSQL
    database to automatically restart and recover following software failures.
    \item \textbf{REL-2:} The Visual Guardian shall implement a configurable
    frame-skip interval (\texttt{PERSON\_DETECTOR\_PROCESS\_EVERY}) to ensure
    sustained CPU stability during continuous 24-hour monitoring.
    \item \textbf{REL-3:} The system shall gracefully handle temporary sensor
    disruptions, such as microphone disconnections or remote inference endpoint
    timeouts, without terminating the monitoring process.
\end{enumerate}

\subsection{Usability}
\begin{enumerate}
    \item \textbf{USE-1:} Alert messages shall be presented in unambiguous natural
    language, requiring no specialised training for nursing staff to interpret.
    \item \textbf{USE-2:} The dashboard shall use a distinct colour-coded scheme
    (Red for Critical, Orange for High) to enable immediate visual differentiation
    of alert severity.
\end{enumerate}

\subsection{Performance}
\begin{enumerate}
    \item \textbf{PER-1:} The end-to-end latency, from the moment a critical event
    occurs to the alert appearing on the dashboard, shall be less than 2 seconds
    under normal operating conditions.
    \item \textbf{PER-2:} The Visual Guardian shall maintain a minimum processing
    rate of 10 frames per second on the target edge hardware under the
    \texttt{LOCAL} inference mode.
    \item \textbf{PER-3:} The Cognitive Core LLM reasoning shall complete Tier-2
    verification in under 500 milliseconds on average.
\end{enumerate}

\subsection{Security}
\begin{enumerate}
    \item \textbf{SEC-1:} All raw video frame processing (YOLOv11n person detection
    and MoViNet action classification) shall be executed locally on the edge device,
    without transmitting raw video feeds to any external server.
    \item \textbf{SEC-2:} The Gemini Cognitive Core shall receive only text-based,
    structured event metadata for Tier-2 and Tier-3 reasoning, and shall never
    receive raw visual or audio streams.
    \item \textbf{SEC-3:} The FastAPI-based dashboard shall restrict access to
    historical event logs stored in the PostgreSQL database, permitting viewing
    only to authorised medical personnel.
\end{enumerate}
